% Chapter 4 -- Rubric: PO3 (3.3, 3.4), PO5 (5.1)
\chapter{System Design and Implementation}
\label{ch:design}

\section{Functional Design and Sub-problem Partitioning}
\label{sec:architecture}
\begin{figure}[htbp]
  \centering
  \figplaceholder{Architecture diagram}
  \caption{System architecture}
  \label{fig:architecture}
\end{figure}
\guidance{Partition the problem into sub-problems or modules; explain each one's
responsibility and how they interact. Show where the standards and codes of
practice from \Cref{sec:standards} and any health, safety, or environmental
considerations shaped the design.}

\section{Design Refinement and Simulation}
\label{sec:detailed-design}
\guidance{Present the detailed design using the views that fit the project (ERD,
class, sequence, DFD, API, wireframes, circuit, model). Include any design
calculations or estimates made (e.g. expected load or storage), the
analyses, prototypes, or simulations used to verify the design against the
requirements, and how the design was refined to facilitate implementation.}

\section{Application of Tools}
\label{sec:tools}
\guidance{Justify the selection of the important tools and technologies and
state their relevant limitations.}

\section{Implementation}
\label{sec:implementation}
\guidance{Describe how the key modules were implemented and the development
process followed (repository, version control, code review, CI/CD).}

\section{Deployment}
\label{sec:deployment}
\guidance{Describe how the system is deployed and operated: the target
environment and infrastructure, configuration and release process,
monitoring and backup, and who operates and maintains it. State the
computing resources the deployment consumes; their environmental impact is
assessed in \Cref{sec:sustainability}.}